%---------------------------------------------------------------------
%	Abstract
%---------------------------------------------------------------------
\chapter{Abstract}				
\label{ch:abs}
%---------------------------------------------------------------------
% TAKEN FROM UCT WEBSITE
% The Doctoral Degrees Board recommends that candidates include an abstract 
% that fits onto one page, and which includes the author's full name, thesis 
% title and date. The text should not exceed 350 words. Candidates may also 
% include a more substantial summary of their work, in addition to the abstract.
% Follow the same format for an MSc. abstract.

% The abstract should stand on its own. It should answer the following questions:
%	1.	What did the author do? What ideas, notions, hypotheses, 
%		concepts, theories or thoughts were investigated?
%	2.	How did the author do the work? What data were generated and used? 
%		What was the origin of the data? How were data gathered? 
%		What tests, scales, indices, or summary measures were used? 
%		In other words, how was the analysis and/or synthesis done?
%	3.	What were the conclusions and what were the significant findings?
%
% Some studies cannot readily be summarised in this way and require more 
% descriptive abstracts. Do not use telegraphic phrases. Do not repeat 
% information given in the title. Do not use abbreviations. The purpose of 
% an abstract is to enable a researcher/examiner to understand the essential 
% hypothesis, method and findings of the research.
%---------------------------------------------------------------------
\begin{center}
	\textbf{\Large \titl}\\
			\vskip 0.2cm
			\auth\\
			\vskip 0.2cm
	\textit{\footnotesize \FullDate}
			\vskip 1cm
\end{center}
%---------------------------------------------------------------------
% Write abstract here [350 words]
%---------------------------------------------------------------------
Antarctic sea ice regulates the Southern Ocean's uptake of heat and carbon dioxide, much of it at the ice underside.
In the marginal ice zone this underside has not yet been reconstructed optically, as satellites cannot observe it
and sonar cannot resolve its centimetre-scale structure. This project investigated whether the stereo cameras of a
BlueROV2 remotely operated vehicle can reconstruct an overhead underwater surface.

A test rig, calibration board and set of targets were built, and three datasets were recorded in a swimming pool.
Printed pyramid and egg carton tiles of known geometry served as ground truth, while painted foam blocks imitated
features of the ice underside.

Merging greyscale and colour channel detections and tuning the detector raised the underwater tag detection rate
from 35 to 43 per cent. Five calibration approaches were compared on repeated held-out splits. A Brown-Conrady
model fitted to underwater images gave a held-out reprojection error of 0.29 pixels, against 8.6 pixels for an
in-air calibration and 1.9 pixels for the Pinax refraction model, and measured tags on an independent target to
within about 2~mm.

Classical and learnt stereo matchers were then compared for depth, along with four camera pose methods and three
ways of fusing frames into a surface. The learnt FoundationStereo matcher gave the most complete depth maps and was
the only matcher to keep the full ice shape in direct sunlight. COLMAP gave the most accurate poses in shade, with
an absolute trajectory error of 6.8~mm. The final pipeline combines COLMAP poses, FoundationStereo depth and
truncated signed distance function fusion, and reconstructed the printed pyramids with median errors of about
2~mm. The learnt model smoothed over partly hidden curved surfaces, so reconstructions of real ice should be checked
against a second method.

The BlueROV2 cameras can therefore reconstruct an overhead underwater surface with millimetre-scale accuracy. The
dataset and pipeline produced here provide a foundation for developing the platform for under ice deployment.
